\documentclass[11pt,a4paper]{ltjsarticle}

% ── Packages ──
\usepackage[margin=2.5cm]{geometry}
\usepackage{luatexja-fontspec}
\setmainjfont{IPAexMincho}
\setsansjfont{IPAexGothic}
\usepackage{amsmath,amssymb}
\usepackage{booktabs}
\usepackage{longtable}
\usepackage{array}
\usepackage{graphicx}
\usepackage[dvipsnames]{xcolor}
\usepackage{hyperref}
\hypersetup{colorlinks=true,linkcolor=blue,citecolor=blue,urlcolor=blue}
\usepackage{enumitem}
\setlist{nosep}
\usepackage{caption}
\usepackage{float}
\usepackage{tikz}
\usetikzlibrary{arrows.meta,positioning,shapes.geometric,fit,calc}

% ── Layout tweaks ──
\setlength{\parindent}{0pt}
\setlength{\parskip}{0.6em}
\renewcommand{\baselinestretch}{1.15}

% ── Counter setup for supplementary ──
\renewcommand{\thefigure}{S\arabic{figure}}
\renewcommand{\thetable}{S\arabic{table}}
\renewcommand{\thesection}{S\arabic{section}}
\renewcommand{\theequation}{S\arabic{equation}}

% ══════════════════════════════════════════════════════════════════
\begin{document}

\begin{center}
{\LARGE\bfseries 補足資料}

\vspace{0.5em}
{\large タンパク質配列からの階層的糖鎖予測により\\
構造クラスレベルでの決定境界が明らかに}

\vspace{0.5em}
{\normalsize 松井 祐介$^{1,2}$}
\end{center}

\vspace{1em}
\tableofcontents
\vspace{1em}

% ══════════════════════════════════════════════════════════════════
\section{ETLパイプラインとナレッジグラフの構築}
\label{sec:etl}

% ── Figure S1 ──
\begin{figure}[H]
\centering
\begin{tikzpicture}[
    node distance=0.6cm and 1.5cm,
    dbbox/.style={draw, rounded corners, fill=#1!15, minimum width=2.4cm,
                  minimum height=0.7cm, font=\small\bfseries, align=center},
    stagebox/.style={draw, rounded corners, fill=gray!20, minimum width=2.2cm,
                     minimum height=0.7cm, font=\small\bfseries},
    arrow/.style={-{Stealth[length=2mm]}, thick},
    label/.style={font=\scriptsize, text=gray!70!black},
]

% Data sources
\node[dbbox=blue]   (glygen)    {GlyGen};
\node[dbbox=green,  below=of glygen]    (glytoucan) {GlyTouCan};
\node[dbbox=orange, below=of glytoucan] (uniprot)   {UniProt};
\node[dbbox=purple, below=of uniprot]   (chembl)    {ChEMBL};
\node[dbbox=brown,  below=of chembl]    (ptmcode)   {PTMCode};
\node[dbbox=cyan,   below=of ptmcode]   (reactome)  {Reactome};

% Stage 1: Download
\node[stagebox, right=2.5cm of glytoucan] (download) {ダウンロード};
\node[label, above=0.1cm of download] {レート制限付きAPI};

% Stage 2: Clean
\node[stagebox, below=of download] (clean) {クリーニング};
\node[label, right=0.1cm of clean, anchor=west] {\begin{tabular}{l}ID正規化\\スキーマ検証\end{tabular}};

% Stage 3: Build
\node[stagebox, below=of clean] (build) {ビルド};
\node[label, right=0.1cm of build, anchor=west] {\begin{tabular}{l}重複排除\\エッジ構築\end{tabular}};

% Stage 4: Validate
\node[stagebox, below=of build] (validate) {バリデーション};
\node[label, right=0.1cm of validate, anchor=west] {\begin{tabular}{l}孤立ノード検出\\自動修正\end{tabular}};

% Output
\node[draw, rounded corners, fill=yellow!15, minimum width=2.5cm,
      minimum height=1.2cm, font=\small, align=center,
      right=2.5cm of clean] (kg) {\textbf{glycoMusubi}\\78,263 nodes\\2,508,960 edges};

% Arrows: sources → download
\foreach \src in {glygen, glytoucan, uniprot, chembl, ptmcode, reactome}
    \draw[arrow, gray] (\src.east) -- (download.west |- \src.east) -- (download.west);

% Arrows: pipeline stages
\draw[arrow] (download) -- (clean);
\draw[arrow] (clean) -- (build);
\draw[arrow] (build) -- (validate);

% Arrow: build → KG
\draw[arrow] (build.east) -- (kg.west);

% Dashed box around pipeline
\node[draw, dashed, gray, rounded corners, inner sep=0.4cm,
      fit=(download)(validate)] {};

\end{tikzpicture}
\caption{\textbf{glycoMusubi ETLパイプラインの概要。}
6つの公開データベースが4段階パイプラインを通じて統合される：(1)~リトライロジック付きレート制限APIクライアントによるダウンロード、(2)~識別子正規化とスキーマ検証によるクリーニング、(3)~重複排除とエッジ構築によるビルド、(4)~孤立ノード検出とオプションの自動修正によるバリデーション。パイプラインは完全に自動化され、単一コマンド（\texttt{python scripts/pipeline.py}）で再現可能である。}
\label{fig:etl_pipeline}
\end{figure}

% ══════════════════════════════════════════════════════════════════
\section{ナレッジグラフスキーマ}
\label{sec:schema}

% ── Table S1: Node types ──
\begin{table}[H]
\centering
\caption{\textbf{glycoMusubiのノード型と識別子フォーマット。}}
\label{tab:node_types}
\small
\begin{tabular}{llllr}
\toprule
\textbf{ノード型} & \textbf{主要ソース} & \textbf{IDフォーマット} & \textbf{例} & \textbf{件数} \\
\midrule
糖鎖    & GlyTouCan  & \texttt{G\textbackslash d\{5\}[A-Z]\{2\}} & G00055MO & 58,628 \\
タンパク質   & UniProt    & Swiss-Prot AC     & P12345   & 4,372 \\
酵素    & GlyGen     & Swiss-Prot AC     & Q10469   & 1,459 \\
部位      & UniProt, PTMCode & \texttt{SITE::ID::pos::res} & SITE::P12345::142::N & 10,218 \\
疾患   & UniProt    & 自由テキスト         & Type 2 diabetes & 1,936 \\
化合物  & ChEMBL     & \texttt{CHEMBL\textbackslash d+} & CHEMBL12345 & 746 \\
バリアント   & UniProt    & 構造化ID     & VAR\_012345 & 522 \\
モチーフ     & GlyTouCan  & 自由テキスト         & Lewis X   & 181 \\
反応  & Reactome   & 自由テキスト         & R-HSA-123 & 164 \\
細胞内局在 & UniProt & \texttt{LOC::name} & LOC::Golgi & 37 \\
\midrule
\multicolumn{4}{l}{\textbf{合計}} & \textbf{78,263} \\
\bottomrule
\end{tabular}
\end{table}

% ── Table S2: Relation types ──
\begin{table}[H]
\centering
\caption{\textbf{glycoMusubiの関係型とエッジ数。}}
\label{tab:relation_types}
\small
\begin{tabular}{lllr}
\toprule
\textbf{関係} & \textbf{ソース型} & \textbf{ターゲット型} & \textbf{エッジ数} \\
\midrule
parent\_of         & 糖鎖   & 糖鎖   & --- \\
child\_of          & 糖鎖   & 糖鎖   & --- \\
subsumes           & 糖鎖   & 糖鎖   & --- \\
subsumed\_by       & 糖鎖   & 糖鎖   & --- \\
has\_motif         & 糖鎖   & モチーフ    & --- \\
produced\_by       & 糖鎖   & 酵素   & --- \\
has\_product       & 反応 & 糖鎖   & --- \\
catalyzed\_by      & 反応 & 酵素   & --- \\
has\_glycan        & タンパク質/酵素 & 糖鎖 & --- \\
has\_site          & タンパク質/酵素 & 部位   & --- \\
glycosylated\_at   & 部位     & 糖鎖   & --- \\
inhibits           & 化合物 & 酵素   & --- \\
associated\_with\_disease & タンパク質/酵素 & 疾患 & --- \\
has\_variant       & タンパク質  & バリアント  & --- \\
ptm\_crosstalk     & 部位     & 部位     & --- \\
localized\_in      & タンパク質/酵素 & 細胞内局在 & --- \\
\midrule
\multicolumn{3}{l}{\textbf{総エッジ数}} & \textbf{2,508,960} \\
\bottomrule
\end{tabular}
\end{table}

注：関係型ごとのエッジ数は糖鎖の構造的階層エッジ（parent\_of、child\_of、subsumes、subsumed\_by）が大部分を占め、250万エッジの大半を構成する。

% ══════════════════════════════════════════════════════════════════
\section{O結合型糖鎖検索結果}
\label{sec:olinked}

% ── Table S3: O-linked retrieval ──
\begin{table}[H]
\centering
\caption{\textbf{O結合型糖鎖検索性能（部位レベル、v6モデル）。}
O結合型糖鎖の高いMRR（0.847）はN結合型糖鎖の低いMRR（0.075）と対照的であり、N結合型予測の課題が本質的な予測不可能性ではなく、より大きく多様な候補空間から生じるという仮説を支持する。}
\label{tab:olinked}
\small
\begin{tabular}{lrrrrrr}
\toprule
\textbf{糖鎖型} & \textbf{MRR} & \textbf{H@1} & \textbf{H@3} & \textbf{H@10} & \textbf{H@50} & \textbf{MR} \\
\midrule
O結合型   & 0.847 & 0.775 & 0.918 & 0.956 & 0.975 & 6.1 \\
N結合型   & 0.075 & 0.026 & 0.067 & 0.182 & 0.391 & 203.3 \\
全体    & 0.314 & 0.257 & 0.330 & 0.421 & 0.572 & 142.4 \\
\midrule
\multicolumn{7}{l}{\textit{人気度ベースライン MRR = 0.285}} \\
\bottomrule
\end{tabular}
\end{table}

v6モデルはN結合型とO結合型の糖鎖を同時に訓練する（訓練サンプル90,735、テストサンプル9,288、総糖鎖候補2,469）。O結合型糖鎖空間はN結合型空間よりも大幅に小さく多様性が低いため、ほぼ完璧な検索が可能である。この結果は検索アーキテクチャ自体が高い能力を持つことを実証している。低いN結合型MRRはモデルの限界ではなく予測問題の困難さの結果である。

% ══════════════════════════════════════════════════════════════════
\section{ナレッジグラフ埋め込みモデルの比較}
\label{sec:kg_embedding}

% ── Table S4: KG embedding baselines ──
\begin{table}[H]
\centering
\caption{\textbf{KG埋め込みモデルの比較（400エポック、232,726テストトリプル）。}
すべてのモデルは標準的なフィルタードランキングプロトコルで評価される。TransEは支配的な階層関係に対する有効性により最高の全体MRRを達成する。GlycoKGNetは困難な\texttt{has\_glycan}関係でより高いMRRを達成する（表~\ref{tab:per_relation}参照）。}
\label{tab:kg_baselines}
\small
\begin{tabular}{lrrrrr}
\toprule
\textbf{モデル} & \textbf{MRR} & \textbf{H@1} & \textbf{H@3} & \textbf{H@10} & \textbf{MR} \\
\midrule
TransE (ベースライン)    & 0.474 & 0.411 & 0.495 & 0.598 & 560 \\
RotatE (ベースライン)    & 0.295 & 0.181 & 0.335 & 0.531 & 407 \\
DistMult (ベースライン)  & 0.042 & 0.017 & 0.038 & 0.082 & 4,009 \\
\midrule
GlycoKGNet R1        & 0.244 & 0.119 & 0.294 & 0.495 & 686 \\
GlycoKGNet R3 ($k=1$)  & 0.232 & 0.121 & 0.275 & 0.444 & 727 \\
GlycoKGNet R3 ($k=32$) & 0.223 & 0.101 & 0.274 & 0.480 & 2,100 \\
\midrule
GlycoKGNet R2 (帰納的) & 0.207 & 0.085 & 0.237 & 0.522 & 374 \\
\bottomrule
\end{tabular}
\end{table}

\textbf{全体MRRと糖鎖特異的性能に関する注記。}
全体MRRは糖鎖階層関係（parent\_of、subsumes）が支配的であり、テストトリプルの大部分を構成する。TransEはこれらの構造関係で優れている。生物学的に情報量の多い\texttt{has\_glycan}関係については、GlycoKGNetがTransEを大幅に上回る（表~\ref{tab:per_relation}）。帰納的GlycoKGNet R2モデルは訓練中に見ていないエンティティを埋め込めるため、すべての下流予測タスクに使用される。

% ══════════════════════════════════════════════════════════════════
\section{関係ごとの性能分析}
\label{sec:per_relation}

% ── Table S5: Per-relation metrics ──
\begin{table}[H]
\centering
\caption{\textbf{TransEベースラインとGlycoKGNetの関係ごとのMRR比較。}
GlycoKGNetは生物学的に重要な\texttt{protein::has\_glycan}関係で改善された性能を示す（3.5倍の改善）。一方TransEは糖鎖の構造的階層関係で優位である。帰納的GlycoKGNet R2は未見エンティティへの強い汎化を維持する。}
\label{tab:per_relation}
\small
\begin{tabular}{lrrr}
\toprule
\textbf{関係} & \textbf{TransE} & \textbf{GlycoKGNet R4} & \textbf{帰納的 R2} \\
 & MRR & MRR & MRR \\
\midrule
glycan::produced\_by::enzyme        & 0.668 & 0.384 & 0.604 \\
reaction::has\_product::glycan      & 0.506 & 0.391 & 0.542 \\
glycan::has\_motif::motif           & 0.337 & ---   & 0.474 \\
glycan::subsumes::glycan            & 0.403 & ---   & 0.286 \\
glycan::parent\_of::glycan          & 0.378 & ---   & --- \\
compound::inhibits::enzyme          & 0.289 & ---   & --- \\
\midrule
protein::has\_glycan::glycan        & 0.028 & 0.096 & 0.024 \\
enzyme::has\_glycan::glycan         & 0.049 & ---   & 0.022 \\
protein::has\_site::site            & 0.000 & ---   & 0.002 \\
\bottomrule
\end{tabular}
\end{table}

\texttt{has\_glycan}関係は埋め込みモデルにとって本質的に困難である。なぜなら各タンパク質は細胞コンテキストに応じて多くの異なる糖鎖にリンクされ得るため、高分散のテール分布を生成するからである。この関係のすべてのモデルにわたる低いMRRは、階層的ベンチマークで観察された決定境界をさらに支持する。

% ══════════════════════════════════════════════════════════════════
\section{階層的予測の詳細結果}
\label{sec:detailed_results}

% ── Table S6: Detailed task results ──
\begin{table}[H]
\centering
\caption{\textbf{すべての階層的予測タスクの詳細結果。}}
\label{tab:detailed_tasks}
\small
\begin{tabular}{llrrrrrr}
\toprule
\textbf{タスク} & \textbf{モデル} & \textbf{F1/MRR} & \textbf{Prec./H@1} & \textbf{Rec./H@3} & \textbf{AUC/H@10} & \textbf{Acc.} & \textbf{MR} \\
\midrule
\multicolumn{8}{l}{\textit{二値分類タスク}} \\
N結合型分類 & ESM-2 MLP & 0.924 & 0.882 & 0.972 & 0.825 & --- & --- \\
O結合型分類 & ESM-2 MLP & 0.641 & 0.593 & 0.698 & 0.716 & --- & --- \\
\midrule
\multicolumn{8}{l}{\textit{ランキング・回帰タスク}} \\
部位ランキング & ESM-2 v5 & 0.688$^*$ & --- & 0.683$^\dagger$ & 0.742 & 0.688 & --- \\
タイプ予測 & Frac. regr. & --- & --- & --- & --- & 0.844 & --- \\
\midrule
\multicolumn{8}{l}{\textit{構造予測タスク}} \\
ファミリー（8クラス） & ESM-2 MLP & --- & --- & --- & --- & 0.638$^a$ & --- \\
クラスタ($K$=20) & Site retr. & 0.397$^*$ & 0.166 & 0.458 & 0.626 & --- & --- \\
\quad クラスタ H@$k$ & & --- & --- & 0.665$^b$ & 0.786$^c$ & 0.978$^d$ & --- \\
\midrule
\multicolumn{8}{l}{\textit{正確な糖鎖検索}} \\
部位レベル & v6b retr. & 0.066$^*$ & 0.025 & 0.060 & 0.145 & --- & 233 \\
タンパク質レベル & v3 retr. & 0.118$^*$ & --- & --- & 0.228 & --- & --- \\
\bottomrule
\end{tabular}

\vspace{0.3em}
\raggedright\scriptsize
$^*$MRR; $^\dagger$Recall@3; $^a$Top-1 accuracy (top-2 = 0.858);
$^b$H@3; $^c$H@5; $^d$H@10.
\end{table}

% ══════════════════════════════════════════════════════════════════
\section{帰納的評価の詳細}
\label{sec:inductive}

% ── Table S7: Inductive evaluation ──
\begin{table}[H]
\centering
\caption{\textbf{GlycoKGNet R2の帰納的評価。}
1,184エンティティ（全エンティティの15\%）がそのすべてのエッジ（63,063トリプル）とともに訓練中にホールドアウトされる。モデルはモダリティ特異的特徴量のみを使用してこれらの未見エンティティを埋め込む必要がある。}
\label{tab:inductive}
\small
\begin{tabular}{lrrrrr}
\toprule
\textbf{設定} & \textbf{MRR} & \textbf{H@1} & \textbf{H@3} & \textbf{H@10} & \textbf{MR} \\
\midrule
全帰納的トリプル & 0.207 & 0.085 & 0.237 & 0.522 & 374 \\
\bottomrule
\end{tabular}
\end{table}

\begin{table}[H]
\centering
\caption{\textbf{関係ごとの帰納的MRR内訳。}
糖鎖の構造的特徴を含む関係（produced\_by、has\_motif、has\_product）は未見エンティティへの汎化が良好である。\texttt{has\_glycan}と\texttt{has\_site}の関係は帰納的設定でも困難のままである。}
\label{tab:inductive_per_relation}
\small
\begin{tabular}{lrrrr}
\toprule
\textbf{関係} & \textbf{MRR} & \textbf{H@1} & \textbf{H@10} & \textbf{MR} \\
\midrule
glycan::produced\_by::enzyme    & 0.604 & 0.558 & 0.694 & 210 \\
reaction::has\_product::glycan  & 0.542 & 0.493 & 0.668 & 175 \\
glycan::has\_motif::motif       & 0.474 & 0.440 & 0.535 & --- \\
glycan::subsumes::glycan        & 0.286 & 0.000 & 0.934 & 178 \\
\midrule
protein::has\_glycan::glycan    & 0.024 & 0.008 & 0.045 & 3,145 \\
enzyme::has\_glycan::glycan     & 0.022 & 0.000 & 0.035 & 328 \\
protein::has\_site::site        & 0.002 & 0.000 & 0.000 & 5,338 \\
\bottomrule
\end{tabular}
\end{table}

% ══════════════════════════════════════════════════════════════════
\section{細胞型酵素発現解析の詳細}
\label{sec:celltype_details}

\subsection{酵素アノテーションカバレッジ}

本研究の重要な発見は酵素-糖鎖アノテーションの疎さである。表~\ref{tab:enzyme_coverage}はこのギャップを定量化する。

\begin{table}[H]
\centering
\caption{\textbf{glycoMusubiにおける酵素-糖鎖アノテーションカバレッジ。}}
\label{tab:enzyme_coverage}
\small
\begin{tabular}{lr}
\toprule
\textbf{統計} & \textbf{値} \\
\midrule
テストセット中のN結合型糖鎖総数      & 2,357 \\
1個以上の酵素アノテーションを持つ糖鎖  & 929 (39.4\%) \\
3個以上の酵素アノテーションを持つ糖鎖 & --- \\
TSアトラス中の糖鎖修飾酵素遺伝子  & 339 \\
Tabula Sapiensの細胞型            & 748 \\
\midrule
酵素カバレッジを持つ構造クラスタ（$K = 20$） & 1 of 20 \\
メガクラスタ（最大）中の糖鎖       & 1,843 (78.2\%) \\
\bottomrule
\end{tabular}
\end{table}

酵素アノテーションが単一のメガクラスタに集中していることは、酵素発現シグナルが20の構造クラスタのうち19を区別できないことを意味する。メガクラスタ内では、1,843の糖鎖が限られた酵素リンクのセットに基づいてランキングを競うため、解像度が低くなる。

\subsection{評価モードの詳細}

4つの評価モードは細胞型情報の使用方法が異なる：

\begin{enumerate}
\item \textbf{周辺}：クラスタ内の各糖鎖$g$について、酵素スコアを全748細胞型にわたって平均する：
$\bar{s}(g) = \frac{1}{|C|} \sum_{c \in C} s(g, c)$。
これは「コンテキストフリー」の酵素シグナルを表す。

\item \textbf{オラクル}：各テストサンプルについて、真の糖鎖のランクを最大化する細胞型$c^*$が選択される：
$c^* = \arg\max_c \text{rank}(g_{\text{true}}, c)$。
これは完璧な細胞型同定の利点の上限を提供する。

\item \textbf{人気度}：糖鎖は訓練セットの出現頻度（正例としての出現回数）によりランキングされる。糖鎖の頻度分布は高度に偏っているため、これは強力なベースラインである。

\item \textbf{学習スコアラー}：InfoNCE損失で訓練されたニューラルネットワークが酵素発現ベクトルと糖鎖埋め込みを共有空間にマッピングする。モデルには頻度効果を捕捉するための糖鎖ごとのバイアス項が含まれる。
\end{enumerate}

% ══════════════════════════════════════════════════════════════════
\section{ESM-2特徴量抽出}
\label{sec:esm2}

タンパク質表現はESM-2（6.5億パラメータ）[本文15]から導出される。残基レベル埋め込み（1,280次元）が最終トランスフォーマー層から抽出され、\texttt{.pt}ファイルとしてキャッシュされる。部位レベルタスクでは、糖鎖付加部位位置での埋め込みが直接使用される。タンパク質レベルタスクでは、全残基にわたる平均プーリングが適用される。

\textbf{カバレッジ。}
残基レベルESM-2埋め込みはKG中の9,138タンパク質のうち1,611（17.6\%）で利用可能である。これはESM-2抽出がGlyConnectで確認された糖鎖付加部位を持つタンパク質についてのみ実施されたためである。ベンチマーク評価のすべてのタンパク質についてESM-2埋め込みが利用可能である。

% ══════════════════════════════════════════════════════════════════
\section{データ分割戦略}
\label{sec:splits}

すべての予測タスクはタンパク質で層化された単一のtrain/validation/testスプリットを使用する：各タンパク質とそのすべての関連糖鎖付加部位は正確に1つのスプリットに出現する。これにより共有タンパク質特徴量を通じた情報漏洩を防止する。

\begin{table}[H]
\centering
\caption{\textbf{N結合型部位レベル検索タスク（v6b）のデータ分割サイズ。}}
\label{tab:splits}
\small
\begin{tabular}{lrr}
\toprule
\textbf{スプリット} & \textbf{サンプル数} & \textbf{割合} \\
\midrule
訓練      & 70,791 & 81.3\% \\
バリデーション &  ---   & --- \\
テスト       &  8,708 & 10.0\% \\
\midrule
総候補数 & 2,357 & N結合型糖鎖 \\
\bottomrule
\end{tabular}
\end{table}

KG埋め込みモデル（GlycoKGNet R2）は別個の帰納的スプリットを使用する：エンティティの15\%がそのすべてのエッジとともにホールドアウトされ、1,184の未見エンティティと63,063の帰納的テストトリプルが提供される。

% ══════════════════════════════════════════════════════════════════
\section{再現性}
\label{sec:reproducibility}

実験パイプライン全体は以下のコマンドで再現可能である：
\begin{verbatim}
python scripts/reproduce_paper.py
\end{verbatim}

このスクリプトは9つのステージを順次実行する：KG構築、KG埋め込み訓練、N結合型分類、部位ランキング、タンパク質レベル糖鎖検索、部位レベル糖鎖検索、細胞型条件付き予測、結果統合および図表生成。\texttt{--dry-run}フラグは実行せずにコマンドを表示し、\texttt{--stage N}フラグはステージ$N$から再開する。

\textbf{ハードウェア要件。}
KG埋め込み訓練には$\geq$8\,GB VRAMのGPUが必要（NVIDIA RTX 3090でテスト済み）。ESM-2特徴量抽出には$\geq$4\,GB VRAMが必要。その他のすべてのタスクはCPUで30分以内に実行可能。

\textbf{ソフトウェア。}
Python 3.11, PyTorch 2.1+, PyTorch Geometric 2.4+。固定された依存関係を持つDockerイメージが提供される。

\end{document}
